%!TEX program = xelatex
\let\stop\empty

\documentclass[12pt,a4paper]{exam-zh}

\usepackage{xeCJK}
\setCJKmainfont{Noto Serif CJK SC}
\usepackage{geometry}
\usepackage{array}
\usepackage{tikz}
\usepackage{amsmath}
\usepackage{wrapfig}
\usepackage{adjustbox}
\usepackage{enumitem}
\usepackage{fancyhdr}
\usepackage{pgfplots}
\pgfplotsset{compat=1.18}
\usetikzlibrary{calc}

\usepackage{scalerel}
\def\mcirc{\mathbin{\scalerel*{\circ}{j}}}
\def\msquare{\mathord{\scalerel*{\Box}{gX}}}
\geometry{margin=2cm}

\AtBeginDocument{\zihao{4}}
% ----------- 调整 exam -----------


%region exam setup
\examsetup{
    % problem/label = \arabic*., 
    % 设置大题题号格式 
        question = {
        top-sep = 1em,      % 题头上方的间距（默认 0.25em）
        bottom-sep = 1em,   % 题号下方的间距（默认 0pt）
    },
    % 隐藏大题后的分值显示 
    problem/show-points = false,
    % 调整小题列表的缩进和间距
    list/step-name = {}, 
    solution = {
    show-solution = show-move, % 答案自动移动到卷末
    label-content = {answer},  % 答案区的标题文字
    show-qed = true            % 在解答末尾显示结束符号 ◻ 
  },
    page = {
    size = a4paper,           % 设定纸张大小 
    foot-content = { - ; - }  % 页脚显示为 -1-, -2- 格式 
  }
}
\setlist[enumerate,1]{
    label=(\alpha*), 
    itemsep=0.5em, 
    leftmargin=2em, 
    after=\hfill % 配合您之前要求的百分比靠右样式
}
\newlist{questionparts}{enumerate}{1}
\setlist[questionparts]{
    label=(\alph*),
    itemsep=0.5em,
    leftmargin=2em,
    rightmargin=4em
}
\newcommand{\qscore}[1]{%
    \unskip\nobreak\hfill
    \makebox[0pt][l]{\makebox[4em][r]{(#1\%)}}%
}
\newcommand{\solutionlayout}{%
    \setlength{\abovedisplayskip}{0.6em}%
    \setlength{\belowdisplayskip}{0.6em}%
    \setlength{\abovedisplayshortskip}{0.35em}%
    \setlength{\belowdisplayshortskip}{0.35em}%
    \setlength{\jot}{0.35em}%
}
%endregion




\begin{document}

%region 范围：高三统计学-集中趋势 | 考点：算术平均数（未分组数据） | 难度：基础
\begin{question}[points = 4]
    某教师记录 5 位学生某次测验的分数如下：
    \[
        76,\ 84,\ 90,\ 68,\ 82
    \]
    求这组数据的平均数。
    \hfill (4\%)
\end{question}

\begin{solution}
    求 5 个数据的平均数。 \\
    直接利用平均数公式 $\bar{x}=\dfrac{\sum x_i}{n}$。
    \begin{align*}
        \bar{x} &= \frac{76+84+90+68+82}{5} \\
        &= \frac{400}{5} \\
        &= 80
    \end{align*}
    \hfill (4分)
\end{solution}
%endregion

%region 范围：高三统计学-集中趋势 | 考点：加权平均数（频数分布） | 难度：中等
\begin{question}[points = 6]
    某班 30 位学生进行一次共 5 题的随堂测验，每人答对题数 $x$（单位：题）的频数分布如下表：

    \begin{center}
    \begin{tabular}{c|ccccc}
        答对题数 $x$ & 1 & 2 & 3 & 4 & 5 \\
        \hline
        人数 $f$ & 2 & 5 & 10 & 8 & 5
    \end{tabular}
    \end{center}

    求该班学生平均答对题数。
    \hfill (6\%)
\end{question}

\begin{solution}
    求频数分布表所示数据的加权平均数。 \\
    利用加权平均数公式 $\bar{x}=\dfrac{\sum f_ix_i}{\sum f_i}$。
    \begin{align*}
        \sum f_i &= 2+5+10+8+5 = 30
    \end{align*}
    \hfill (2分)
    \begin{align*}
        \sum f_ix_i &= 1(2)+2(5)+3(10)+4(8)+5(5) \\
        &= 2+10+30+32+25 \\
        &= 99
    \end{align*}
    \hfill (2分)
    \[
        \bar{x} = \frac{99}{30} = 3.3
    \]
    \hfill (2分)
\end{solution}
%endregion

%region 范围：高三统计学-集中趋势 | 考点：中位数（未分组数据，无需公式） | 难度：基础
\begin{question}[points = 4]
    某教师记录 10 位学生完成一份作业所需的时间（单位：分钟）如下：
    \[
        23,\ 19,\ 31,\ 27,\ 15,\ 35,\ 29,\ 41,\ 18,\ 25
    \]
    求这组数据的中位数。
    \hfill (4\%)
\end{question}

\begin{solution}
    求 10 个数据的中位数。 \\
    将数据由小到大排列，取中间两项的平均数。
    \[
        15,\ 18,\ 19,\ 23,\ 25,\ 27,\ 29,\ 31,\ 35,\ 41
    \]
    \hfill (2分)
    \begin{align*}
        M &= \frac{25+27}{2} \\
        &= 26
    \end{align*}
    \hfill (2分)
\end{solution}
%endregion

%region 范围：高三统计学-集中趋势 | 考点：中位数（分组数据，插值公式） | 难度：中等
\begin{question}[points = 8]
    某次数学测验中，50 名学生的成绩分布如下表：

    \begin{center}
    \begin{tabular}{|c|c|c|c|c|c|}
        \hline
        成绩 & 50-60 & 60-70 & 70-80 & 80-90 & 90-100 \\
        \hline
        人数 $f$ & 5 & 12 & 20 & 9 & 4 \\
        \hline
    \end{tabular}
    \end{center}

    求这组数据的中位数。
    \hfill (8\%)
\end{question}

\begin{solution}
    求分组数据的中位数。 \\
    先确定中位数所在组，再代入中位数插值公式 $M=L+\left(\dfrac{\frac{n}{2}-F}{f_m}\right)c$。
    \begin{align*}
        n &= 5+12+20+9+4 = 50, \quad \frac{n}{2}=25
    \end{align*}
    \hfill (2分)
    \begin{align*}
        F(\text{累积至}\ 70\text{-}80\ \text{之前}) &= 5+12 = 17
    \end{align*}
    \hfill (2分)
    中位数组为 $70\text{-}80$，即 $L=70,\ F=17,\ f_m=20,\ c=10$。
    \begin{align*}
        M &= 70 + \left(\frac{25-17}{20}\right)(10) \\
        &= 70 + \frac{80}{20} \\
        &= 70+4 \\
        &= 74
    \end{align*}
    \hfill (4分)
\end{solution}
%endregion

%region 范围：高三统计学-集中趋势 | 考点：众数（分组数据，插值公式） | 难度：中等
\begin{question}[points = 8]
    某校体育测试记录 70 名学生一分钟跳绳个数，分布如下表：

    \begin{center}
    \begin{tabular}{|c|c|c|c|c|c|}
        \hline
        跳绳个数 & 50-60 & 60-70 & 70-80 & 80-90 & 90-100 \\
        \hline
        人数 $f$ & 8 & 12 & 24 & 16 & 10 \\
        \hline
    \end{tabular}
    \end{center}

    求这组数据的众数。
    \hfill (8\%)
\end{question}

\begin{solution}
    求分组数据的众数。 \\
    确定众数组，代入众数公式 $Mo=L+\left(\dfrac{d_1}{d_1+d_2}\right)c$，其中 $d_1=f_m-f_1,\ d_2=f_m-f_2$。
    众数组为频数最高的 $70\text{-}80$，即 $L=70,\ f_m=24,\ f_1=12,\ f_2=16,\ c=10$。
    \begin{align*}
        d_1 &= 24-12 = 12 \\
        d_2 &= 24-16 = 8
    \end{align*}
    \hfill (3分)
    \begin{align*}
        Mo &= 70 + \left(\frac{12}{12+8}\right)(10) \\
        &= 70 + \left(\frac{12}{20}\right)(10) \\
        &= 70+6 \\
        &= 76
    \end{align*}
    \hfill (5分)
\end{solution}
%endregion

\end{document}